\begin{table}[htbp]
\centering
\small
\caption{Уральская панель (9 регионов): эффект не идентифицируется}
\label{tab:ural}
\begin{tabular}{lrrrrr}
\toprule
Зависимая & Спецификация & Коэффициент & Ст. ошибка & N & Кластеров \\
\midrule
Общие выдачи & базовая & +0.064 & (0.071) & 216 & 9 \\
Общие выдачи & + контроли × пост & -0.046 & (0.134) & 216 & 9 \\
Общие выдачи & + тренды + контроли & -0.078 & (0.257) & 216 & 9 \\
Рыночные & базовая & +0.250 & (0.236) & 216 & 9 \\
Рыночные & + контроли × пост & +0.228 & (0.182) & 216 & 9 \\
Рыночные & + тренды + контроли & +1.249*** & (0.279) & 216 & 9 \\
Льготные без массовой & базовая & +2.059* & (1.127) & 214 & 9 \\
Льготные без массовой & + контроли × пост & +0.965** & (0.479) & 214 & 9 \\
Льготные без массовой & + тренды + контроли & +3.698*** & (0.232) & 214 & 9 \\
\bottomrule
\end{tabular}
\par\smallskip\begin{minipage}{0.92\textwidth}\footnotesize
Те же зависимые переменные, что и в основной таблице, но выборка ограничена девятью субъектами Уральского макрорегиона. Контроли — население, доля детей 0--6 лет, структура семей по ВПН-2020, цена первичного рынка 2023 года. Знак коэффициента меняется от спецификации к спецификации; рандомизационный тест даёт $p=0{,}655$, плацебо-событие значимо.
\end{minipage}
\end{table}
